\section{QueryConditionedF and Heuristic Pre-Training}
\label{sec:query-conditioned-f}

\subsection{Free-energy functional}

For state $\rho$ and query embedding $q \in \mathbb{R}^{384}$:
\begin{equation}
F(\rho, q) = \underbrace{\sum_{s} D_{\mathrm{KL}}(\rho_s \,\|\, \pi_{\mathrm{prior},s})}_{\mathrm{complexity}}
         + \underbrace{\tfrac{1}{2\sigma^2} \| q - g_{\mathrm{JEPA}}(\rho) \|^2}_{\mathrm{accuracy\;(JEPA)}}
         + \underbrace{\lambda_J \sum_{s,t} J_{s,t} \langle \rho_s, \rho_t \rangle}_{\mathrm{coupling}}
\label{eq:F}
\end{equation}

The decoder $g_{\mathrm{JEPA}} : \Delta^{128} \to \mathbb{R}^{384}$
is a two-layer MLP (128 $\to$ 256 $\to$ 384, GELU activations)
with parameters pre-trained as described below. The prior
$\pi_{\mathrm{prior}, s}$ defaults to the uniform simplex; in our
experiments we keep $\sigma^2 = 1$ and $\lambda_J = 0.1$. The
species coupling matrix $J \in \mathbb{R}^{4 \times 4}$ is taken
from the Levelt-Baddeley specification used by
\texttt{T2FreeEnergy} in our codebase.

\paragraph{Choice of $\sigma^2$ and $\lambda_J$.}
The default $\sigma^2 = 1$ normalizes the JEPA accuracy to the
same scale as the KL complexity (both in log-units, given our
embedding normalization). $\lambda_J = 0.1$ was fixed by a
coarse pilot (three values $\{0.01, 0.1, 1.0\}$ on a 1\,k-query
subset): $0.01$ underweights the species coupling, producing
species drift decoupled from Levelt-Baddeley expectations;
$1.0$ over-regularizes, collapsing the state toward the
coupling null-space. We do not claim $0.1$ is optimal;
sensitivity is deferred to the companion full paper, where it
is tested as a free hyperparameter under joint
$f_\theta$/$g_{\mathrm{JEPA}}$ training.

\subsection{Gradients}

The complexity term yields
$\partial_{\rho_s} = \log(\rho_s / \pi_{\mathrm{prior},s}) + 1$
and the coupling term yields
$\partial_{\rho_s} = 2 \lambda_J \sum_{t} J_{s,t} \rho_t$,
both analytical. The JEPA term is differentiated via JAX autodiff.
One JKO step minimizes
$F(\rho, q) + \tfrac{1}{2h} W_2^2(\rho, \rho_\tau)$,
the Wasserstein formulation of variational inference. As
$\tau \to \tau+1$ the flow approximates the active-inference
posterior $q^\star(s \mid q)$.

\subsection{Heuristic labeler (phase A)}
\label{sec:labeler}

To pre-train $g_{\mathrm{JEPA}}$ without ground-truth species
labels, we build a deterministic French NLP labeler that maps each
query $q$ to a target
$\pi^{\mathrm{lbl}} \in (\Delta^{32})^4$
(\cref{alg:labeler}):

\begin{itemize}
    \item \textbf{phono}: 32 phonetic classes (4 broad groups
        $\times$ 8 finer subclasses) extracted via the
        \texttt{espeak-ng} backend of \texttt{phonemizer}.
    \item \textbf{sem}: 32 categorical semantic bins, currently
        keyed via a stable lemma hash (a placeholder for a
        WordNet-fr or KMeans-MiniLM mapping; this affects only
        the absolute identity of bins, not the simplex
        structure).
    \item \textbf{lex}: 32 log-frequency bins computed from
        Lexique.org 3.83~\citep{newboucard2004lexique}.
    \item \textbf{syntax}: 32 Universal Dependency labels,
        counted from a SpaCy-fr dependency
        parse~\citep{honnibal2020spacy}.
\end{itemize}

\begin{algorithm}[ht]
  \caption{Heuristic labeler (phase A target)}
  \label{alg:labeler}
  \small
  \begin{tabbing}
    ~~~\=\quad\=\quad\=\quad\kill
    \textbf{Input:} query string $q$ \\
    \textbf{Output:} target distribution
        $\pi^{\mathrm{lbl}} \in (\Delta^{32})^4$ \\[1ex]
    \> $\pi_{\mathrm{phono}} \leftarrow$ normalize(count\_phoneme\_classes(phonemizer($q$))) \\
    \> $\pi_{\mathrm{sem}}   \leftarrow$ one\_hot(stable\_hash(spacy\_lemmas($q$)) mod 32) \\
    \> $\pi_{\mathrm{lex}}   \leftarrow$ normalize(count\_freq\_bins(lexique\_lookup(tokens($q$)), 32)) \\
    \> $\pi_{\mathrm{syntax}} \leftarrow$ normalize(count\_ud\_labels(spacy\_parse($q$), 32)) \\
    \> \textbf{return} $(\pi_{\mathrm{phono}}, \pi_{\mathrm{sem}}, \pi_{\mathrm{lex}}, \pi_{\mathrm{syntax}})$
  \end{tabbing}
\end{algorithm}

Phase A then minimizes
$\| f_\theta(q) - g_{\mathrm{JEPA}}(\mathrm{flatten}(\pi^{\mathrm{lbl}})) \|^2$
via Adam-W on $\sim$5--10\,k French queries; $f_\theta$ is held
fixed during phase A and selected by the encoder ablation reported
in Section~\ref{sec:experiments}.

\paragraph{Why hash-MLP wins despite a smaller parameter count.}
The Section~\ref{sec:experiments} ablation finds that a
$0.65$\,M-parameter hash-MLP encoder (architecture C) outperforms a
$2.44$\,M-parameter distilled-MiniLM analogue (B) on the pilot and
scale phases. We hypothesize this is not a defect of the larger
model but an architectural match: the 4 $\times$ 32 simplex target
has limited effective capacity (a 128-d simplex is isomorphic to a
$127$-d manifold in $\mathbb{R}^{128}$), and the hash-MLP produces
representations with sparsity structure that aligns better with the
per-species histograms of Algorithm~\ref{alg:labeler}. The
MiniLM analogue embeds semantic content that is partially wasted at
the simplex resolution. A companion full paper (in preparation)
tests this hypothesis directly by ablating the
$g_{\mathrm{JEPA}}$ hidden dimension against each encoder family.

\subsection{Ablation}
\label{sec:ablation}

Figure~\ref{fig:ablation} reports the per-species KL on the FR
hybrid corpus at 10\,k (pilot) and 50\,k (scale) for three text
encoders ($\mathrm{B}_{\text{distilled}}$,
$\mathrm{C}_{\text{hash-mlp}}$,
$\mathrm{D}_{\text{tiny-tf}}$).
Figure~\ref{fig:scaling} plots the total test KL trajectory
$10\text{\,k} \to 50\text{\,k}$ per architecture.
$\mathrm{C}_{\text{hash-mlp}}$ retains rank-1 at both scales and
the gap to rank-2 $\mathrm{B}_{\text{distilled}}$ collapses from
$\approx 3\,\%$ at 10\,k to $\approx 0.3\,\%$ at 50\,k, confirming
the R3 kill-switch promotion to Top-2.

\begin{figure}[t]
  \centering
  \includegraphics[width=\linewidth]{figures/ablation_kl_species.pdf}
  \caption{Per-species KL across encoders on pilot10k and scale50k.}
  \label{fig:ablation}
\end{figure}

\begin{figure}[t]
  \centering
  \includegraphics[width=0.75\linewidth]{figures/scaling_curves.pdf}
  \caption{Total test KL scaling $10\text{\,k} \to 50\text{\,k}$
    per architecture.}
  \label{fig:scaling}
\end{figure}
